\chapter{CƠ SỞ LÝ THUYẾT} \label{chap:background}

This chapter presents a structured review of key techniques that enhance quadrotor safety and resilience during flight. It is organized into three main sections. The first section introduces collision avoidance strategies based on vision and perception, focusing on how drones detect and react to nearby obstacles in advance. The second section explores advanced control methods, with an emphasis on the $\mathcal{L}_1$Quad approach, which enables the quadrotor to maintain stability and adapt to external disturbances such as collisions or deformation of Tombo propellers. The final section discusses sensor-based collision detection, highlighting methods that utilize onboard sensors and passive mechanical responses to identify and localize collision events in real time.

\section{Collision Avoidance strategies}
Numerous collision avoidance strategies have been developed for \gls{UAV}s to ensure safe navigation, especially in cluttered or dynamic environments. These methods typically fall into categories such as geometric planning, force-field approaches, optimization-based schemes, and sense-and-avoid systems~\cite{avoidancetypes,avoidancetypes2}. Central to these strategies is integrating onboard sensors that perceive the environment and guide the \gls{UAV}’s response.

Vision-based sensors are widely adopted due to their rich spatial sensing capabilities and low payload weight. Among these, monocular vision offers the simplest hardware setup, relying on a single camera to extract visual features from sequential images. While machine learning methods such as YOLO-based detectors or neural networks~\cite{YoloFastest} can improve detection accuracy, they often require significant onboard computational power. Simpler alternatives include image processing approaches that identify color or shape-defined objects~\cite{monocolor}, but they struggle with unknown or unstructured obstacles. 

In~\cite{monocularcamera}, monocular vision was used for tasks such as landing on moving platforms. However, it remains insufficient for detecting large or static obstacles like buildings. To address this, stereo vision systems are employed to generate disparity maps for depth estimation. For instance,~\cite{avoidancetypes} developed a biologically inspired stereo system mounted on a servo, allowing owl-like head movements to track dynamic obstacles. Another approach, "pushbroom stereo," combines stereo imagery with forward motion and lightweight processing for fast obstacle detection~\cite{pushbroom}.